
%\section{Proposed Method: SAGE}
%\label{sec:method}

We propose SAGE (\textbf{S}ubject-\textbf{A}gnostic to \textbf{G}eometrically-aligned \textbf{E}mbeddings), a two-stage framework designed to bridge the intra-/cross-subject gap in EEG-to-image decoding.  

\section{Subject-Disentangled Training}
% Here we will merge encoder + features + InfoNCE + bottleneck + SubjectMix BUT FORCUS ON JUBJECT MIX  
\label{sec:training}
\paragraph{Problem Formulation}
%\label{sec:problem}
We consider a multi-subject EEG-to-image decoding setting. Let $\xv_{s,i} \in \R^{C\times T}$   denote the EEG signal recorded from subject $s \in \mathcal{S}$ in response to image $i \in \mathcal{I}$,  where $C=63$ is the number of EEG channels and $T=250$ the number of time steps.  
Let $\vv_i \in \R^{d_v}$ denote the corresponding image embedding.   
We observe a dataset $\Dc = \{(\xv_{s,i}, \vv_i): s\in \Sc, i \in \Ic\}$ where multiple subjects provide distinct EEG responses to the same visual stimulus.

To avoid relying on explicit subject identification, we model EEG signals as inherently noisy observations arising from both stimulus-dependent and subject-specific factors. Specifically, we assume:
\begin{equation} \label{eq:genmodel}
  \xv_{s,i} = g(\zv_i, \uv_s) + \epsilonv_{s,i}, % yet introduce noise here ... keep simple model 
\end{equation}
 where $\zv_i \in \Zc$ is a latent variable associated with the stimulus and shared across subjects, $\uv_s \in \Uc$ captures subject-specific variability, and $\epsilonv_{s,i}$ represents intrinsic noise in EEG measurements, including physiological variability and sensor noise.

Our objective is to learn an encoder $E_\theta : \R^{C \times T} \to \R^d$ that aligns EEG responses with their corresponding image embeddings while removing subject-specific variability. 
\begin{comment}{\color{red}Under this model, we aim to recover a representation that depends only on the stimulus factor $\zv_i$. 
In particular, we seek an encoder such that there exists a function $h : \Zc \to \R^d$ satisfying
\begin{equation}
  E_\theta(\xv_{s,i}) = h(\zv_i),
\end{equation}
for all subjects $s$ and stimuli $i$.}
\end{comment}
This induces a natural invariance structure: for a fixed stimulus $i$, the set $\{ \xv_{s,i} : s \in \Sc \}$ corresponds to different realizations of the same latent factor. 
In practice, this amounts to enforcing consistency across subject-conditioned observations associated with the same image, while maintaining discriminability across different stimuli.

In the following, we introduce a training strategy that explicitly promotes this invariance.



 
 

We now describe how we learn EEG representations that align with image embeddings while suppressing subject-specific variability.

\paragraph{SubjectMix: enforcing invariance.}
To enforce the invariance structure introduced above, we propose \emph{SubjectMix}, a cross-subject mixing regularization. 

For each stimulus $i$, let define the set of original EEG samples
\begin{equation}
\Xc_i = \{ \xv_{s,i} : s \in \Sc \}.
\end{equation}
We construct mixed samples by combining pairs of subjects, and define the corresponding index set
\begin{equation}
\widetilde{\Sc} = \{ (s,s') : s \neq s',\; s,s' \in \Sc \},
\end{equation}
with associated samples
\begin{equation}
\widetilde{\Xc}_i = \{ \xv_{s,s',i} = \Ac(\xv_{s,i}, \xv_{s',i}) : (s,s') \in \widetilde{\Sc} \},
\end{equation}
where $\Ac: \R^{C \times T} \times \R^{C \times T} \rightarrow \R^{C \times T}$ is an aggregation operator designed to preserve the stimulus-dependent factor while attenuating subject-specific variability.
 
Under the observation model \eqref{eq:genmodel}, such transformations primarily perturb subject-specific variability while preserving the stimulus-dependent factor $\zv_i$.
We therefore enforce
\begin{equation}
     E_\theta(\xv_{k,i}) \approx E_\theta(\xv_{k',i}), \quad \forall\, k \in \Sc,\; k' \in  \widetilde{\Sc},
\end{equation}
encouraging representations that depend only on $\zv_i$.

\paragraph{Representation learning.}
We model the encoder as a composition of temporal and spatial operators followed by a projection head: Divide $ \theta = (\theta_p, \theta_s, \theta_t)$, we model 
\begin{equation}
E_\theta = P_{\theta_p} \circ S_{\theta_s} \circ T_{\theta_t},
\end{equation}
where $T_{\theta_t}$ captures temporal dynamics along the time axis, $S_{\theta_s}$ aggregates information across EEG channels, and $P_{\theta_p}$ projects the resulting features into a shared embedding space.

Given an input $\xv_{s,i} \in \R^{C \times T}$, with $s\in \Sc \cup \widetilde{\Sc}$ and $i \in \Ic$, the embedding is
\begin{equation}
\ev_{s,i} = E_\theta(\xv_{s,i}) \in \R^d.
\end{equation}

The operator $S_{\theta_s}$ acts across all $C=63$ channels, allowing the model to integrate distributed cortical activity, in contrast to approaches restricted to a subset of electrodes~\cite{du2025shallow}.
The projection head $P_{\theta_p}$ consists of a lightweight MLP with residual connection and normalization.

We further constrain the representation to lie in a compact space of dimension $d=64$, which acts as a structural bottleneck. 
This restriction limits the capacity to encode subject-specific variability while preserving stimulus-dependent information, thereby supporting the invariance objective.
 
\paragraph{Training objective.}
%{\color{red}Let $\ev_{s,i} = E_\theta(\xv_{s,i}) \in \R^d$ denote the (unnormalized) EEG embedding of subject $s$ viewing stimulus $i$ , and let $\hat{\vv}_i \in \R^d$ denote the $\ell_2$-normalized image embedding  EST CE QUE CEST NECESSAIRE ? CITE PAPER}. 
We use a multi-positive contrastive objective based on stimulus-grouped batches~\cite{khosla2020supcon}. Each batch contains $M$ distinct stimuli, each associated with $K$ EEG responses from different subjects, yielding $MK$ EEG samples. 
%For a given EEG embedding $\ev_{s,i}$,  all embeddings $\vv_{i}$ corresponding to the same stimulus form the positive set, while embeddings of other stimuli act as negatives.
{\color{red}For a given EEG embedding $\ev_{s,i}$, the corresponding image embedding $\vv_{i}$ is treated as the positive target, while image embeddings $\{\vv_{j}\}_{j \neq i}$ from other stimuli act as negatives.}

The EEG-to-image loss is defined as:
\begin{equation}
\mathcal{L}_{\mathrm{E \to V}} =
-\frac{1}{M} \sum_{i=1}^{M} \frac{1}{K}\sum_{s=1}^{K}
\log 
\frac{\exp(\ev_{s,i} \cdot \vv_{i} / \tau)}
{\sum_{j=1}^{M} \exp(\ev_{s,i} \cdot \vv_{j} / \tau)}.
\end{equation}

Symmetrically, the image-to-EEG loss considers all EEG embeddings associated with the same stimulus as positives:
\begin{equation}
\mathcal{L}_{\mathrm{V \to E}} =
-\frac{1}{M} \sum_{i=1}^{M}
\frac{1}{K} \sum_{s=1}^{K}
\log 
\frac{\exp(\vv_{i} \cdot \ev_{s,i} / \tau)}
{\sum_{(s',j)} \exp(\vv_{i} \cdot \ev_{s',j} / \tau)}.
\end{equation}

The final contrastive loss is
\begin{equation}
\mathcal{L}_{\mathrm{ctr}} = \frac{1}{2} \left( \mathcal{L}_{\mathrm{E \to V}} + \mathcal{L}_{\mathrm{V \to E}} \right).
\end{equation}





\section{Transductive Geometric Alignment}
\label{sec:testing}
% merge ZCA + CSLS + Sinkhorn + Procrustes + iterative loop and  genralization of rotaation 
Despite subject-disentangled training, a residual geometric mismatch remains between the EEG embedding space of a new subject and the image embedding space. 
While the training objective encourages invariance, it does not uniquely fix the representation, which remains defined up to subject-dependent transformations. 
These transformations arise from factors such as cortical geometry and electrode placement, leading to systematic distortions across subjects.

To address this, we adopt a transductive approach: given a set of $N$ EEG queries from an unseen subject and $N$ candidate image embeddings with a known one-to-one structure, we estimate a global alignment without requiring labels. 
We treat the two modalities as point clouds and iteratively estimate correspondences and a global transformation.

\paragraph{Alignment pipeline.}
At test time, we consider a single unseen subject and omit the subject index for simplicity. 
Let

$$\Em = [\ev_1^\top,\ldots,\ev_N^\top]^\top \in \mathbb{R}^{N\times d},
\qquad
\Vm = [\vv_1^\top,\ldots,\vv_N^\top]^\top \in \mathbb{R}^{N\times d}$$

denote the EEG and image embedding matrices, respectively. 
The correspondence between EEG and image embeddings is unknown. 
We therefore align the two sets through the following steps.


\textbf{(i) Whitening.}
We first normalize the EEG embeddings using the empirical mean and covariance of the test EEG set:
\begin{equation}
\muv = \frac{1}{N}\sum_{i=1}^{N} \ev_i,
\qquad
\Sigmam = \frac{1}{N}\sum_{i=1}^{N}(\ev_i-\muv)(\ev_i-\muv)^\top .
\end{equation}
To improve numerical stability, we use a shrinkage covariance estimate
\begin{equation}
\widetilde{\Sigmam}
=
(1-\rho)\Sigmam
+
\rho \frac{\operatorname{tr}(\Sigmam)}{d}\I_d .
\end{equation}
The whitened and normalized EEG embeddings are
\begin{equation}
\widetilde{\ev}_i
=
\frac{
\widetilde{\Sigmam}^{-1/2}(\ev_i-\muv)
}{
\left\|\widetilde{\Sigmam}^{-1/2}(\ev_i-\muv)\right\|
}.
\end{equation}
We denote by $\widetilde{\Em}\in \mathbb{R}^{N\times d}$ the matrix whose rows are $\widetilde{\ev}_i$.

\textbf{(ii) Similarity correction.}
We compute a corrected similarity matrix $\Sm \in \mathbb{R}^{N\times N}$ using CSLS:
\begin{equation}
s_{i,j}
=
\mathrm{CSLS}(\widetilde{\ev}_i,\vv_j)
=
2\cos(\widetilde{\ev}_i,\vv_j)
-
r_E(\widetilde{\ev}_i)
-
r_V(\vv_j),
\end{equation}
where $r_E$ and $r_V$ denote the average similarity to the nearest neighbors in the opposite embedding set.

\textbf{(iii) Soft assignment.}
We convert the similarity matrix into a soft correspondence matrix using Sinkhorn normalization:
\begin{equation}
\Pm = \mathrm{Sinkhorn}\left(\exp(\Sm/\tau_{\mathrm{sk}})\right),
\end{equation}
where $p_{i,j}$ represents the soft assignment between EEG embedding $\widetilde{\ev}_i$ and image embedding $\vv_j$.

\textbf{(iv) Geometric alignment.}
From the soft assignment matrix $\Pm$, we estimate an orthogonal transformation $\Rm \in \mathbb{R}^{d\times d}$ by solving
\begin{equation}
\Rm^*
=
\arg\min_{\Rm^\top \Rm = \I}
\sum_{i=1}^{N}\sum_{j=1}^{N}
p_{i,j}\left\|\widetilde{\ev}_i\Rm - \vv_j\right\|^2 .
\end{equation}
The optimal solution is given in closed form by $\Rm^* = \Um \Wm^\top,$ where $\Um \Dm \Wm^\top$ is the singular value decomposition of $\widetilde{\Em}^\top \Pm \Vm$.
 

\paragraph{Generalization across samples.}
To assess whether the learned transformation captures intrinsic subject-specific structure rather than overfitting, we estimate the alignment on a subset of test samples and apply it to the remaining ones. 
The resulting improvement demonstrates that the transformation generalizes beyond observed pairs, indicating that it captures consistent geometric structure induced by the subject.










